\section{Discussion}
\label{sec:discussion}

\subsection{Answer to Research Questions}~\label{sec:disc:rq}

We answer the research questions as follows:

\begin{itemize}
  \item \textbf{RQ1:} Our parallel congruence closure algorithm achieves a
  meaningful speedup on random benchmarks, a synthetic benchmark suite designed
  to stress test our algorithm, and a set of circuit equivalence benchmarks. We
  scale well up to 32 cores (\as{give a number here}), but began to plateau
  beyond that, especially on smaller benchmarks.
  \item \textbf{RQ2:} Our algorithm tends to perform better on benchmarks with a
  large number of merges per round. More rounds of congruence closure does not
  have a major effect on parallel speedup.
  \item \textbf{RQ3:} Our algorithm provides a meaningful speedup over a
  sequential implementation on the circuit equivalence benchmarks. It is
  difficult to evaluate our algorithm on SMT or equality saturation benchmarks,
  as congruence closure is part of a larger tool and thus difficult to isolate.
  We leave it to future work to integrate our algorithm into a state-of-the-art
  SMT solver or equality saturation engine.
\end{itemize}

\subsection{Algorithm Modifications and Future Work}~\label{sec:disc:limit-alg}

There are several extensions to the algorithm as presented to be addressed in 
future work.

\textbf{Synchronous Round Structure.} Our synchronous round structure could be
limited in terms of parallel scaling. In particular, if we have congruence
closure problems with large \emph{depth}, but small \emph{width}, our
parallization strategy cannot achieve a good speedup. We implemented an
asynchronous version of the dirty children algorithm, where we did semisort and
unions happened at the same time. Certain threads would semisort and discover
new congruences which the would add to a concurrent bag
datastructure~\cite{concurrentbag}. Then the threads that wer unioning would
probe this concurrent bag for new merges to perform.

We found that this was slower than the synchronous version. A challenge to
address is that even this asynchronous algorithm cannot efficiently parallize a
chain of congruences, i.e. where we you have add the merges $x \equiv y$, $f(x)
\equiv f(y)$, $f(f(x)) \equiv f(f(y))$, etc.



% inherently
% limits parallel scaling. In particular, the instance size processed at each
% round is small relative to the size of the problem for several benchmarks we
% have examined (like the \texttt{egg} benchmarks). We will likely have to take a
% more asynchronous approach to achieve better thread scaling; this would entail
% generating new congruences from discovered equalities immediately instead of
% waiting for the round to complete.

\textbf{Backtracking.} 
\as{Zak can you fill in how we could make our implementation backtrackable}

\textbf{Proof Production.} Our implementation does not produce \emph{proofs}(i.e.
\emph{witnesses}) of the equalities it discovers. For instance if our algorithm
discovers that $x = y$, it cannot replay the chain of equalities and congruences
that led to this conclusion. This is important for a number of reasons: it is
necessary to justify correctness and it is used to produce conflict clauses in
SMT solvers. Our implementation could easily be extended to support proof
production by avoiding path compression in the union-find data structure.
Another option is to maintain two union-find data structures: one for the actual
congruence closure algorithm, and another that maintains the proof forest as
done by the Z3 SMT solver does~\cite{z3-internals}.

\textbf{Incremental Update.} Our algorithms could support incremental update
(where equalities and new terms may arrive dynamically)

Finally, we would like to support our
implementation to incremental updates (where equalities and new E-node may
arrive dynamically). This is important for both equality saturation and SMT
solvers, where new equalities may be discovered via for example quantifier
instantiation.
